\documentclass[./examen.tex]{subfiles}

\begin{document}
\info{Escuela 4-124 Reynaldo Merín  - Evaluación Electrotécnia II}{TEMAS: f.e.m generada por un conductor que se mueve en un campo magnético, Fuerza producida por un conductor que se mueve en un campo magnético, Ley de Faraday-Lenz, regla de la mano izquierda y derecha.}
\begin{center}
\makebox[\textwidth]{Nombre y Apellido:\enspace\hrulefill}
\end{center}
\begin{questions}
 
 \question Un electroimán tiene una Fuerza de $F=100Kp$ y su inducción magnética en el hierro es de $B=10000Gs$.
 \begin{parts}
    \part[1] Calcular la sección del electroimán $S$ en $cm^2$.
 \vspace{2cm}
    \part[1] Calcular la intensidad de campo magnético $H$ si la permeabilidad relativa del núcleo de electroimán es $\mu_r=2000$
  \vspace{2cm}
 \end{parts}
 
\question El circuito de la figura está construído de chapa magnética $\mu_r=2000$, se desea una inducción de $B=1,8T$ la bobina es de $N=200$ espiras.  
\begin{parts}
\part[1] Calcular la longitud de la línea media $l_m$ en metros. 
\part[1] Calcular la Intensidad de campo en el núcleo $H$ en $Av/m$.
\part[1] Fuerza magnetomotriz $\mathscr{F}$ amperios-vuelta necesarios.
\part[1] Intensidad de corriente que circula por la bobina.
\end{parts}

  
\begin{center}
\begin{tikzpicture}[scale=0.02]
\draw [ultra thick] (0,0) rectangle (180,180);
\draw [ultra thick] (40,40) rectangle (140,140);
\draw [thin,dashed] (20,20) rectangle (160,160);
\draw [thin,Circle -] (-20,140)--(40,140);
\draw [thin] (0,130)--(40,130);
\draw [thin] (0,120)--(40,120);
\draw [thin] (0,110)--(40,110);
\draw [thin] (0,100)--(40,100);
\draw [thin] (0,90)--(40,90);
\draw [thin] (0,80)--(40,80);
\draw [thin] (0,70)--(40,70);
\draw [thin] (0,60)--(40,60);
\draw [thin] (0,50)--(40,50);
\draw [thin,Circle -] (-20,40)--(0,40);
\draw [stealth-stealth] (140,80)--(180,80);
\draw [stealth-stealth] (40,100)--(140,100);
\draw [stealth-stealth] (100,140)--(100,40);
\draw [stealth-stealth] (80,40)--(80,0);

\draw (80,110) node[]{100};
\draw (110,70) node[]{100};
\draw (170,90) node[]{40};
\draw (70,30) node[]{40};

\end{tikzpicture}
\end{center}

\question  Un conductor $l=20cm$ de longitud se mueve dentro de un campo magnético de inducción $B=1Tesla$. 
\begin{parts}
\part[1] Si ejerce sobre el conductor una fuerza de  $F=50N$. ¿Qué intensidad $I$ debe circular por el conductor?
\vspace{1cm}
\part[1] Si la corriente fuese de $I=10A$ ¿Cuánta fuerza $F$ se ejerce sobre el conductor en Newtons?
\vspace{1cm}
\end{parts}
  \question Un solenoide de $N=1000$ espiras, longitud de $l=60cm$, $r=4cm$ y núcleo de madera $\mu_r=1$ el cual tiene arrollado sobre él una bobina de $N=1500$ espiras. 
  \begin{parts}
  \part[1] Calcular la f.e.m. inducida en la bobina cuando se anula en $\Delta t=0,1s$. La corriente es de $I=10A$. 
  \vspace{2cm}
  \end{parts}
  \question[1] Explicar la regla de la mano derecha.
  \end{questions}
 \end{document}